%==============================================================================
\section{Life of a simulation}\label{sec:lifecycle}
%==============================================================================

A steady-state run goes through four objects and one call, always in this order:
\begin{enumerate}[leftmargin=1.5em]
\item \code{ArrayHydroGrid(Nx, Ny, xlims, ylims)} --- the grid (\file{common/grid.jl}). With data on cell centres
  \code{xc}, \code{yc}: \code{xlims, ylims = compute\_lims(xc, yc)}.
\item \code{HydroState(grid, mask, h, b)} --- copies of the geometry, and $N = W = 0$ (\file{common/state.jl}).
\item The model: \code{KazmierczakHydroModel(grid, kappa, abs\_v\_b, A\_visc, G, q\_T; kwargs...)} (all SI, all
  \code{Nx}\,$\times$\,\code{Ny}), or \code{HABHydroModel(grid)}. The K24 constructor checks sizes, converts every number to the
  grid's float type, resolves the option defaults that follow the routing scheme, rejects invalid combinations, and
  allocates all 38 workspace arrays.
\item \code{SteadyStateSimulation(model, grid, state)} (\file{common/simulation.jl}).
\item \code{run!(sim)} = \code{update\_steady\_state!(model, grid, state)}; the results are \code{state.N}, \code{state.W} and the
  model's fields (\code{model.q}, \code{model.Po}, \code{model.N\_inf}, \code{model.psi\_out}, \ldots).
\end{enumerate}
There is no time loop: calling \code{run!} again recomputes everything from the present inputs. A coupled ice model
updates the inputs in place between calls and calls \code{run!} once per step (Section~\ref{sec:coupling}).

\subsection{A minimal script}\label{life:example}
A 30\,km by 30\,km slab with a sloping bed and a floating strip, the K24 defaults plus the regularised Coulomb law:
{\small
\begin{verbatim}
using FastHydrology

Nx, Ny, dx = 30, 30, 1000.0
grid = ArrayHydroGrid(Nx, Ny, (0.0, Nx * dx), (0.0, Ny * dx))
x = [(i - 0.5) * dx for i in 1:Nx, j in 1:Ny]

mask = ones(Nx, Ny); mask[end-2:end, :] .= 0          # floating strip: the grounding line
h = 900.0 .- 0.015 .* x                                 # [m]
b = -150.0 .- 4e-3 .* x                                 # [m]
state = HydroState(grid, mask, h, b)

kappa   = zeros(Nx, Ny)                                 # hard bed
abs_v_b = fill(perYear2perSecond(100.0), Nx, Ny)        # 100 m/yr -> m/s
A_visc  = fill(2.4e-24, Nx, Ny)                         # Pa^-3 s^-1
G, q_T  = fill(0.07, Nx, Ny), fill(0.03, Nx, Ny)        # W/m^2

model = KazmierczakHydroModel(grid, kappa, abs_v_b, A_visc, G, q_T;
            coupling_length_kamb86 = 5.0,               # 5 ice thicknesses
            sliding_law = RegularizedCoulombSlidingLaw(c_till = 0.3))
run!(SteadyStateSimulation(model, grid, state))

state.N, state.W, model.q                               # Pa, m, m^2/s
\end{verbatim}}
A flat bed and surface would give $\nabla\phi_0 = 0$ everywhere: no flow direction, $c_f = 0$, and nothing to route. Real
and test geometries always slope.

\subsection{One steady-state solve}\label{life:solve}
The call tree of \code{update\_steady\_state!} for the default K24 configuration (Warner routing, priority-flood
filling, taped routing, face-averaged $q$, face dissipation on, no friction), with the file of each function:
{\small
\begin{verbatim}
update_steady_state!(model, grid, state)                          run.jl
 |- update_q!(model, grid, state)                                 water_flux.jl
 |   |- update_phi0!                     phi0 = rho_i g h + rho_w g b
 |   |- prepare_routing!(..., ::Warner)                           routing.jl
 |   |   |- coupling_kernel + convolve!  smooth phi0 -> phi0_filled (if lcw > 0)
 |   |   |- potential_filling!(from_true_potential = false)
 |   |   |   '- fill_potential!(::PriorityFloodFill)  heap_push!/heap_pop!
 |   |   |- update_potential_gradients!  |grad phi0| (true), -grad phi0_filled
 |   |   |- copy to minus_grad_phi0_sx/sy, abs_grad_phi0_s
 |   |   '- compute_routing_weights!     w8[d,i,j] via cell_weights!(::Warner)
 |   |- corfac = (|sx| dy + |sy| dx)/|s|
 |   |- invalidate_routing_tape!
 |   '- resolve_q!(..., ::DissipationMeltOn, ::NoFrictionSlidingLaw)
 |       |- update_tau_b!                tau_b = 0                  sliding_law.jl
 |       '- Picard sweeps (<= max_dissipation_iters = 20):
 |           |- mdot_total = mdot_fixed + i_eb
 |           |- add_friction_term!       Q_b = tau_b |u_b|
 |           |- add_dissipation_term!(::FaceDissipation)   from last sweep's faces
 |           |- route_psi_out!(::TapedPsiOut) -> update_psi_out_taped!
 |           |   |- record_routing_tape!  (first sweep only; ..._weights_kernel!)
 |           |   '- replay_routing_tape!
 |           |- update_q_from_psi_out!(model, grid, state)         routing.jl
 |           |   |- update_face_fluxes!  Fx, Fy, diss
 |           |   '- q_from_psi_out!(::QFromFaceAverage)
 |           '- stop when max|q - q_prev| <= 1e-12 max|q| (grounded cells)
 |- update_N!(model, grid, state)                                 effective_pressure.jl
 |   '- update_N_grounded_kernel!   Q, S_inf, H, Po, N_inf, N per cell (threads)
 '- update_W!(model, grid, state)    -> update_W_darcy_weisbach!(::MeanGradient)
\end{verbatim}}
Three variants change the tree:
\begin{itemize}[leftmargin=1.5em]
\item \code{routing\_scheme = GDSWarner()}: \code{prepare\_routing!} fills $\phi_0$ first (\code{JacobiFill}), takes the gradients
  and smooths them (\code{update\_smoothed\_potential\_gradients!}); weights are computed on the fly by
  \code{routing\_weight} and recorded by \code{record\_routing\_tape\_kernel!}; $q = \psiout/c_f$.
\item \code{dissipation\_melt = false} with an $N$-independent law: \code{resolve\_q!} is a single pass, no loop.
\item An $N$-dependent sliding law: \code{resolve\_q!} becomes the joint $(q, N)$ loop, which calls \code{update\_tau\_b!} from the
  current $N$ and \code{update\_N!} inside every sweep (Section~\ref{sec:fixedpoints}). The \code{update\_N!} and \code{update\_W!}
  after \code{update\_q!} then recompute the converged values once more, which is redundant but keeps the call sequence
  the same for every law.
\end{itemize}
$W$ must come after $N$, because \code{ArealConduitThickness} reads $\Sinf$, which \code{update\_N!} computes; nothing reads $W$.

\subsection{HAB}
\code{update\_steady\_state!(::HABHydroModel, ...)} is \code{update\_N!}, which calls \code{update\_Po!} and \code{update\_p\_w!} and then
$N = \max(p_o - p_w, \sigma_t p_o)$, three broadcasts. It reads only \code{state.h}, \code{state.b} and \code{state.mask}.

\subsection{Shakti}
\code{TimeSimulation(ShaktiHydroModel(shakti\_sim))}: \code{run!(sim; kwargs...)} calls \code{Shakti.run!(shakti\_sim; kwargs...)};
\code{step!(sim)} calls \code{Shakti.step!} and advances \code{shakti\_sim.total\_time[]} by \code{dt[]}. The Shakti simulation is
built with Shakti's own API (see the Shakti guide); FastHydrology adds nothing to it.

\subsection{Timings on real grids}\label{life:timing}
The tour's last PART times full solves on the real inputs (one CPU thread for the routing, four for \code{update\_N!}).
On albedo (\code{smp} node) the default K24 solve takes the times below. A compiled solve of the 2\,km Thwaites grid takes
24\,ms. Of that, preparing the routing surface (smoothing, filling, gradients and fractions, done once per solve) is the
largest share at 10\,ms; recording the routing tape costs 4\,ms, also once per solve, and each further Picard sweep then
replays it in 0.5\,ms. The first call costs several times more, because of compilation.
\begin{center}\small
\begin{tabular}{>{\raggedright\arraybackslash}p{8.6cm}>{\raggedleft\arraybackslash}p{2.4cm}}
\toprule
Thwaites 2\,km ($345\times288$, 51\,945 grounded cells), default K24 & time \\
\midrule
first solve, including compilation & 170\,ms \\
full solve, compiled, warm start (Warner, priority flood, taped) & 23.8\,ms \\
\quad \code{update\_phi0!} & $<0.1$\,ms \\
\quad \code{prepare\_routing!}: smoothing, filling, gradients, fractions & 9.8\,ms \\
\quad record the tape and replay it once (154\,885 operations) & 4.1\,ms \\
\quad replay only (one Picard sweep's routing) & 0.5\,ms \\
\quad $\psiout\to q$ (\code{QFromFaceAverage}) & 1.4\,ms \\
\quad \code{update\_N!} (4 threads) & 3.5\,ms \\
\quad \code{update\_W!} & 2.0\,ms \\
full solve, GDSWarner, \code{RecursivePsiOut} & 23.5\,ms \\
full solve, GDSWarner, \code{TapedPsiOut} & 19.9\,ms \\
\midrule
Antarctica 16\,km from a Yelmo restart ($381\times381$, 48\,937 grounded cells), no smoothing & \\
\midrule
first solve / compiled & 48.7\,ms / 24.8\,ms \\
\bottomrule
\end{tabular}
\end{center}
Times are from the archived run \file{test/code\_tour/logs/code\_tour\_final\_48451698.out} (\code{sbatch} on an albedo \code{smp} node, \code{julia -t 4}); each compiled time is the best of three after one untimed
call. A cold solve on Thwaites (a fresh model, $q = 0$) needs 8 dissipation sweeps; a warm one starts from the previous $q$.
The recursion did not overflow the stack on Thwaites with GDSWarner in this run, but its depth depends on the flow paths and
the stack limit, so \code{TapedPsiOut} stays the safe choice.

